\chapter{How Imperialism Redrew the World}
\section{A Map Drawn with a Ruler}
Pick up an object: a map.

Not a satellite image or navigation app, but Africa's political map in any current geography textbook. Study its borders for ten seconds. Many are \textbf{perfectly straight}, running hundreds of kilometers as though drawn with a ruler. Egypt--Libya, Libya--Chad, Namibia--Botswana: lines cross deserts, grasslands, and nomadic routes used for hundreds of generations without concern.

Nature has no straight lines. Rivers bend, mountains rise and fall, and the edges of grazing, marriage, and market-going territories blur and merge. Show the continents to someone unfamiliar with world maps and ask which borders were drawn with an office ruler: they would probably choose Africa, correctly.

Compare Europe's jagged boundaries, shaped inch by inch through centuries of war, marriage, and negotiation. Each bend may reflect a river, castle, or battle. Africa's straight lines often follow longitude, latitude, or simply connect two points. Most drawers had never visited and did not need to.

Who redrew the world, when, and with what? What gave them that power? What happened to the people enclosed by the lines? This chapter answers the map's questions.

\section{A Conference Room in Berlin}
Enter a scene.

The afternoon of November 15, 1884, in a high-ceilinged room at German Chancellor Bismarck's Berlin residence. A strong fire burns; cigar smoke hangs in the air. An enormous African map covers the long table, though its interior is mostly blank: coasts, a few explorers' routes, and vast areas marked unknown.

Diplomats from more than a dozen countries sit around it: Britain, France, Germany, Belgium, Portugal, Spain, Italy, Austria-Hungary, Russia, the United States. Formally dressed, speaking diplomatic French, they politely address one another as Your Excellency.

Pause over one fact: \textbf{not one African attended} this discussion of Africa's fate. No king, chief, merchant, or herder was invited; nobody even thought to invite them. A continent of over a hundred million people and centuries-old kingdoms and city-states became blank paper awaiting division.

Meeting until the following February, they produced the Berlin General Act. Among many provisions, two core points stand out: open the Congo basin to every country's free trade; and require future coastal claimants to notify other signatories and establish effective occupation. A drawn circle was insufficient: troops, officials, and flags had to make possession real.

Supposedly cooling rivalry by rules, it did the opposite. Effective occupation fired a starting gun: move first and land is yours. Expeditions, armies, and chartered companies flooded inland, presenting chiefs with foreign-language protection treaties they often could not read. A signature or mark licensed claims to hundreds of thousands of square kilometers. Historians later estimated European control at around a tenth of Africa near 1875 and over nine-tenths near 1900. In a quarter-century, a continent changed owners.

Remember firelight, cigar smoke, and the invisible ruler hovering above that table. Everything ahead answers questions hidden in that room.

\section{A Few Decades, Nine-Tenths of the Land}
State the conflict before answering.

First establish events, the relatively undisputed evidentiary layer. Before the nineteenth century, European colonial footholds largely hugged coasts: ports, forts, trading posts. The final three decades exploded outward. By World War I in 1914, some historians estimate over four-fifths of Earth's land was ruled by Euro-American powers or states founded by their settler descendants. Africa was almost partitioned; India belonged to Britain's Crown. Southeast Asia left Siam, now Thailand, barely independent while France held Vietnam, the Netherlands Indonesia, Britain Burma and Malaya, and Spain then America the Philippines. The Ottoman Empire survived with finances and tariffs controlled by European creditors. China retained government but not secure control of tariffs, concessions, or coastlines.

Now begins the continuing argument over why.

Participants offered one explanation. Berlin's diplomats sincerely saw respectable work: ending slave trading, spreading Christianity, bringing law and railways. Decades later many European textbooks still described imperfect but ultimately civilizing triumph over ignorance.

Those across the map remembered otherwise. Indian railways arrived alongside repeated late-nineteenth-century famines killing millions while grain continued abroad under contract. Leopold II's personally controlled Congo inflicted systematic village violence for rubber quotas, with population losses estimated in millions. Exact numbers remain disputed, not the catastrophe. Chinese treaty-port residents remembered these agreements through national humiliation.

The chapter's real questions emerge in a chain:

What enabled Europe---weapons, money, organization, something else? This is \textbf{capacity}.

What made them want expansion---greed, strategy, belief, competitive compulsion? This is \textbf{motivation}.

Why those decades rather than earlier or later? This is \textbf{timing}.

Most deeply, should we scrutinize how overwhelming power finds reasons for its actions? This is \textbf{justification}.

Before continuing, choose intuitively: was the deepest driver economics, technology, or ideas?

\section[Empire and the People Across the Map]{Empire's Account and the People Across the Map}
Give empire's reasoning its full form. Turning opponents into flat villains is where lazy historical understanding begins. \textbf{Explaining a justification fully is not defending it}. It asks why so many found it persuasive and whether it could return in another costume.

Four layers recur. \textbf{Civilizing mission}: Europeans widely ranked humanity by civilization, placing themselves first and obliged to bring science, law, education, and often Christianity to backward peoples as adults instruct children. \textbf{Free trade as justice}: opening every market benefits all, so trade barriers oppose progress. \textbf{Order and honor}: governments must defend overseas merchants; refusing to compete while others expand means weakness and decline. \textbf{Humanitarianism}: abolition. After centuries participating in Atlantic slavery, Europeans made ending inland African slave trading their moral intervention banner.

Separately, these are not simply madness. Each hides \textbf{who benefits} and assumes the other party has no voice. Railways genuinely arrive, but run from mines to ports for extraction. Trade becomes free, but gunboats negotiate tariffs. Slave trading is condemned while replacement regimes demand forced rubber labor. Understanding this rhetoric's operation is more useful than merely calling it hypocrisy.

Turn the map over and listen to many voices, from very different continental circumstances.

\textbf{South Asia}. India's first conqueror was not Britain's government but a company: the London-listed East India Company, with an army and minting powers, swallowed princely states from the mid-eighteenth century. Only after suppressing the soldier-led 1857 uprising across northern India did government assume company powers. In 1876 Victoria added Empress of India. One country became another crown's jewel.

\textbf{Southeast Asia}. No single power dominated: France took Vietnam, the Dutch Indonesia, Britain Burma and Malaya, and Spain lost the Philippines to America. Java's Dutch Cultivation System required peasants to devote land and time to official export crops: sugarcane, coffee, indigo. Spice and sugar wealth went to Amsterdam while Javanese farmers starved on fertile land.

\textbf{Africa}. We have seen Congo's catastrophe. British West Africa also used chiefs to collect taxes and labor for colonial government: indirect rule reduced costs while converting traditional authority into imperial machinery.

\textbf{West Asia}. The Ottoman Empire avoided formal partition but had European banking supervision of finances and lost independent tariff control: empire without occupation. After its World War I defeat, Britain and France redrew maps through wartime secret agreements, notably Sykes--Picot in 1916, then administered mandates. Some modern Middle Eastern borders were determined at those negotiating tables.

\textbf{East Asia}. China experienced semi-colonialism: a surviving court but treaty losses of tariff, port, concession, and inland navigation powers. Japan opened under gunboat threat, rapidly learned its challengers' methods, then annexed Korea and Taiwan and became imperial. This matters: imperialism is not simply whites ruling people of color but an expansionary logic \textbf{available to anyone who learns it}.

Together, these show imperialism's varied faces: company rule, Crown administration, chiefly intermediaries, financial supervision, treaty ports. The next section classifies them.

\section{Thinking Bridge: Three Forms of Imperial Rule}
Imperialism is easily imagined as one thing: governors, garrisons, flags. Historians studying how the past happened and was written have identified several models, and \textbf{the flagless one can be equally deadly}. Ask of every foreign power which control it exercises.

\textbf{First: direct rule.} The imperial power sends governors and bureaucracies and governs colonies through its own law. Post-1876 India, French Vietnam, and Belgian Congo exemplify it. The highest official is appointed far away; final fiscal and judicial authority is external.

\textbf{Second: indirect rule.} Kings, chiefs, or sultans retain thrones while taxes, armies, and foreign relations are controlled by a neighboring British resident. Britain used this in northern Nigeria; Lugard later explained its cheap, stable use of familiar authority. Local rulers remain, but foreign advisers stand behind their orders. Traditional shells survive with reconstructed interiors.

\textbf{Third: informal empire.} No governor, perhaps no troops beyond a few gunboats, and formal independence survives. Yet treaties surrender key sovereignty: negotiated tariffs, foreign customs management, and exemption of foreign nationals from local courts, or extraterritoriality. Nineteenth-century China and the Ottomans are textbook cases. Inspect treaties, not flags. Who grips fiscal and judicial throats?

This immediately corrects \textbf{the false claim that absence of formal colonization means absence of imperialism}. No single power swallowed China whole, but Shanghai park notices, Robert Hart's nearly half-century leading customs, and treaty tariffs denying protection for domestic industry were real controls. Historians called this the imperialism of free trade. A gun need not fire daily: its harbor presence tilts the negotiating table.

Conversely, these categories overlap. Britain ruled India directly, Nigeria indirectly, and China informally. Methods are toolbox wrenches; the purpose remains making land serve distant interests.

When news mentions influence, spheres of interest, or debt arrangements, try the classification. Which control? Who holds its switch?

\section{A Poem and the Fight Around It}
Read a short primary source.

In 1899, Rudyard Kipling published The White Man's Burden, opening:

\begin{quote}
Take up the White Man's burden—
\end{quote}
The poem urges America, then taking the Philippines, to accept the supposedly thankless but noble duty of civilizing others.

This records empire's self-understanding. Notice its psychology: conquest becomes \textbf{burden}, profit becomes \textbf{sacrifice}---all our hardship is for your good.

The reaction is historically more revealing than the poem. Widely reprinted, it also provoked replies: American anti-imperialists mocked a burden containing others' sugar and land; Indian journalists remarked how profitable carrying it seemed. Over decades the solemn slogan became a sarcastic expression, usually wielded today in mocking quotation marks.

This small source teaches a basic skill: \textbf{who speaks, where, and to whom matters as much as literal wording}. A salesperson's offer to help differs from a neighbor's. Kipling reassures imperial citizens, not colonial subjects. His mobilization addresses rulers' \textbf{self-doubt}: how to hold stolen possessions with an easy conscience.

That unlocks colonial annual reports, mission newsletters, and companies' cheerful shareholder messages. They state facts but first \textbf{account to their own people}. This does not make every word false. Watch both what they tell you and what they need you to believe.

\section{Why Europe? Four Explanations Compete}
Return to why. At least four rival accounts explain Europeans acquiring most of the world within decades. Each illuminates something and leaves blind spots, rather than being simply right or wrong.

\textbf{First: guns, steamships, and quinine, the technological account.}

Industrial Europe produced breech-loading rifles, Maxim guns, shallow-draft steam gunboats, and telegraphs unavailable to earlier colonizers. Firepower gaps became unprecedented. At Omdurman in 1898, Anglo-Egyptian machine guns and artillery defeated tens of thousands of Mahdists with very low own casualties. Steamships penetrated rivers; antimalarial quinine greatly improved European survival in tropical West Africa. The chronology is strong: conquest accelerated as these matured. But technology explains \textbf{could}, not \textbf{wanted}. A sharp knife need not stab. Nor does it explain advanced weapons losing battles.

Consider the essential counterexample, \textbf{Ethiopia}. Menelik II's forces defeated invading Italians at Adwa in 1896, preserving independence. The invaded were not doomed natives awaiting discovery: Ethiopia had monarchy, army, diplomacy, and European-purchased weapons. European invincibility is itself part of victors' narrative.

\textbf{Second: capital seeks outlets, the economic account.}

Early-twentieth-century Hobson offered an argument Lenin developed in Imperialism, the Highest Stage of Capitalism. Industry produces excess goods and capital, domestic markets cannot absorb them, returns fall, and capital needs overseas raw materials, markets, and investment. Colonies are its safety valve. This shifts attention from bad individuals to systemic logic and explains taxpayers funding empire despite no personal colonial interest. But calculations suggest many colonies cost metropolitan society more in military and administrative expenses than they returned, while particular firms and sectors profited: public costs, private profits. And markets alone could use informal control as in China; why expensively rule African deserts?

\textbf{Third: a white racial ranking, the racist-ideology account.}

Nineteenth-century European intellectual life spread pseudoscientific human hierarchies with whites first. This licensed conquest: natural rankings made leaders ruling followers as blameless as gravity. Kipling supplied its literary version. It helps explain violence's \textbf{cruelty}: treating others fully as people would make decades of rubber atrocities difficult. Yet racism often lubricates afterward rather than drives initially; Europeans were brutal to Europeans too. Some racists even opposed imperialism because supposedly inferior peoples were not worth conquest's cost. Ideas following interests may fit better than ideas driving them.

\textbf{Fourth: deny rivals possession, the interstate-competition account.}

Late-nineteenth-century Europe was a powder keg of suspicious powers. A colony's value mattered less than \textbf{keeping it from rivals}, who could use it for coal stations, soldiers, or bargaining. French acquisition demanded a neighboring British claim; latecoming Germany grabbed harder to secure a place in the sun. Berlin convened because rivalry risked shooting. This explains the 1880s acceleration after German unification and American rise reshuffled power and intensified distrust. But distrust itself rests on interests. Without the other motives, insecurity supplies a powder keg without powder.

Combine them for a stronger account: technology gives \textbf{capacity}, capital \textbf{pull}, racism \textbf{permission}, and competition \textbf{the starting gun}. Remove one and events differ. This is not compromise for its own sake: a supposedly complete single explanation has probably discarded something early.

\section{Deeper Challenge: Core, Semiperiphery, Periphery}
Now connect the whole map through Immanuel Wallerstein's 1970s \textbf{world-systems theory}.

Change scale, he says. Asking separately why Britain is rich and India poor asks wrongly. From roughly the sixteenth century, \textbf{one} integrated economic web gradually incorporated the world. It has three positions:

\begin{itemize}
\item \textbf{Core}: controls high-profit manufacturing, finance, and technology and sets rules. Nineteenth-century western Europe.
\item \textbf{Periphery}: compelled to supply cheap materials and labor in an exploited position. Colonies.
\item \textbf{Semiperiphery}: exploited by the core while exploiting places further out, such as Russia then or certain later Latin American countries.
\end{itemize}
The crucial claim is \textbf{inequality between core and periphery results from one system, not two independent development histories}. Richness and poverty are two sides of one coin, like a pump raising one end by draining the other. Ask not why India failed to invent industry but who assigned its role in the web. Eighteenth-century Indian, especially Bengali, cotton conquered European markets; British cloth could not compete. The nineteenth reversed it: India became Britain's major cloth market and raw-cotton supplier while handweaving vanished, a process economic historians call deindustrialization. The people and land did not become less hardworking. Their assigned position changed, through tariffs and guns.

The theory cuts sharply but has limits. Critics say its system resembles a machine with its own will, predetermining everything and reducing struggle, choice, and reversals like Adwa to tiny gear adjustments. Why could Japan move from peripheral incorporation to core within decades when others could not? System-decides-everything and national-character-decides-everything share the mistake of discarding human action.

Its indispensable insight remains: \textbf{empire's departure does not dissolve the web}. That bridges to what follows.

\section[How the Ruled Lived]{Accommodation, Obstruction, Sparks: How the Ruled Lived}
We have examined imperial action; now take the ruled seriously or repeat Berlin's mistake of treating them as blank paper.

Their responses fall roughly into three kinds, often coexisting within one person.

\textbf{First: collaboration.} Today almost an insult, it needs its historical balance of power. Indian princes, northern Nigerian chiefs, and thousands of local colonial employees often chose not resistance versus surrender but how their people could survive, or how new power could restrain old rivals. Calculation, necessity, and genuine advantage coexisted. Decades of rule depended on incorporating locals, not a few thousand Europeans alone. This does not excuse empire: \textbf{recognizing collaborators as reasoning human beings reveals that imperial rule relied on division and incorporation, not merely guns}.

\textbf{Second: adaptation.} Learn the ruler's language, schools, examinations, and tools. Meiji Japan adapted nationally; young Indians, Vietnamese, and West Africans learning English or French adapted personally. Ironically, colonial education often produced nationalism's first organizers, who turned learned nation, freedom, and rights against empire's legitimacy. Tools offered by the master dismantled the master's house.

\textbf{Third: resistance.} Some was loud: India's 1857 uprising, African anti-French and anti-British wars, China's Boxers and later revolutions. Some was quiet and long ignored: slow work, feigned ignorance, tax evasion, noncooperation, indigenous gods smuggled into the ruler's religion. Studying Southeast Asian peasants, James Scott called these weapons of the weak. Direct confrontation meant death; small daily jams accumulated into change. \textbf{Silence is not consent; apparent obedience may conceal sophisticated passive resistance}.

Then sparks became national movements. The 1885 Indian National Congress initially petitioned within the system; decades later Gandhi turned nonviolent noncooperation into mass spinning, salt making, and boycotting British cloth. Mid-twentieth-century independence swept Asia and Africa: India in 1947, Indonesia and Vietnam through war, African flags rising throughout the 1950s and 1960s. In 1955, more than twenty Asian and African states met at Bandung, Indonesia: the ruled entered their own conference room as hosts for the first time, seventy years after Berlin.

\section{Ghosts on the Map}
Berlin's meeting ended in 1885, but its remains lie under today's footsteps. Identify continuing colonial legacies, distinguishing established facts, reasonable explanations, and open questions.

\textbf{Borders are facts.} Many African segments still follow coordinates or straight connections, splitting ethnic groups and joining historical adversaries, one background to numerous sub-Saharan conflicts. Some Middle Eastern borders trace postwar Anglo-French negotiations. These are not conflicts' sole causes---blaming every war on a century ago is lazy---but denying any connection is untenable.

\textbf{Languages are facts.} Independent India could not choose one indigenous national language without privileging its speakers, so colonial English remained as a compromise, now a ladder into global technology. Francophone West Africa and Anglophone East Africa still shape textbooks and overseas study. Your own hours learning English are among empire's deepest legacies.

\textbf{Economic structures combine fact and explanation.} Many former colonies still depend on one or two raw exports---cocoa, coffee, oil, copper---their assigned roles. How much continuation follows colonial paths, such as port-directed railways and clerical education, and how much post-independence governance? World-systems theory emphasizes the former, critics the latter. Real money rides on this debate because answers about poverty determine whether rich countries owe something.

\textbf{Memory and dignity remain active questions.} European museums hold colonial loot amid intermittent return negotiations. Some cities remove colonial statues while others defend them. Textbooks continually dispute how to tell colonial history: understatement becomes forgetting, emphasis accusations of hatred. Everyday campus echoes ask why English scores matter so much, why international standards often default to Europe and America, and why English brand names seem prestigious. These feelings are the conference room's century-long shadow.

Fairness requires adding that legacies are not only shackles. Railways, modern schools, and administrative frameworks also help independent states build: Indian railways now carry Indians. Yet routes, curricula, and administrative languages retain original intentions. This leads to the final question.

\section[Settling an Old Account]{Your Question: How Should an Old Account Be Settled?}
Return to the opening map.

You know who drew the straight lines, with what power and justifications, what people on both sides experienced, and how lines extend into English textbooks and news feeds.

Your question is: \textbf{how should we settle colonialism's century-old account of borders, languages, railways, institutions, debts, artifacts, and wounds today?}

Give an answer and reasons, weighing both sides first.

Those demanding reckoning have a case: stolen objects should return, extracted wealth deserves compensation, and ignoring continuing inequality compounds the debt. If waiting a century cancels injury, what remains of justice?

Those looking forward have one too: perpetrators and victims have died, so is charging today's taxpayers for great-grandparents fair? Does blaming history excuse current rulers? Would destroying railways, schools, and shared English as wholly negative assets hurt poor people most?

You also know the harder fact: no balance sheet can contain Congo's dead, India's starved, or divided communities. Uncalculable does not mean unaccountable, just as inability to repay does not cancel obligation. One legacy can be both shackle and tool, depending on present use.

Should an uncalculable old debt be written off, repaid in installments, or calculated differently? What determines whether a history should be remembered?

There is no standard answer. But next time you see a map, restitution dispute, or history told differently by countries, you already know the ruler exists. Seeing it begins every answer.
